\section{从一支流动作战军到许多可供养的军府}
\label{sec:regional-polities-and-military-fiscal-order}

广明元年十二月，黄巢军入长安，建立大齐；僖宗逃往西南，次年在成都维持行在。两件事合在一起，才显出晚唐崩解的实际尺度：都城可以被攻取，皇帝的诏令仍可以在另一座城市发出，而决定战争能否延续的，已是哪些军队能够得到粮食、道路、仓储与赏赐。不能把黄巢军只写成一股从关东卷向京城的破坏力量，也不能把成都行在写成完整无缺的中央政府。前者进入长安后要面对关中供给和官僚组织，后者则依赖西川的城防、粮运与地方军政资源；两边都在试图把临时的军事聚合变成可以持续供养的秩序。%
\footnote{黄巢入长安及建大齐见账本 LT-880-062，僖宗至成都并维持行在见 LT-881-063；基本年序可据《旧唐书·僖宗本纪》《旧唐书·黄巢传》《新唐书·黄巢传》及《资治通鉴》唐纪相关条。入城秩序、居民态度与成都行在的实际行政覆盖均缺少连续记录，不能用本纪的朝廷视角补足。}

这一变化也改写了将领与朝廷相遇的方式。中和二年，朱温自黄巢军降唐，获授宣武军节度使；次年，唐军与沙陀、河东等军收复长安；黄巢于884年兵败后，其部众、武器和行军经验没有随首领死亡而消失，而被不同军镇和地方武装吸收。这里的关键不是把朱温、李克用或沙陀军预先安排成后来王朝的主角，而是看见朝廷为收复都城所作的授官、赏赐和借兵，怎样同时增加了若干军府在兵员、声望与征发上的资本。收复长安恢复的是朝廷的名义中心，并没有把军队重新放回一条由京师独自供给、独自指挥的链条。%
\footnote{朱温降唐见 LT-882-064，收复长安见 LT-883-065，黄巢败亡及余部去向的限度见 LT-884-066。《旧唐书·朱全忠传》《旧唐书·李克用传》《旧唐书·黄巢传》与《资治通鉴》保存了基本序列；它们分别服务于唐朝复国、后梁建国或后唐军功的叙事，不能据此精算各军人数，或断言所有齐军余部都转入同一军镇。}

盐池收入引起的885年关中危机，把这种军事化的财政关系露在了表面。田令孜意图调动河中军并控制盐利，王重荣拒命，僖宗又一次离开长安。盐并非抽象的“国库收入”：谁能取得盐利，谁便较有能力支付军队、维持本镇的文书与人事，并在京畿冲突中同皇帝和宦官讨价还价。材料不足以把每一笔盐利还原成账目，更不能据一场冲突推算晚唐全国财政；但它足以说明，朝廷、宦官与节度使争夺的并非单纯的官位，而是使兵权能够延续的收入节点。%
\footnote{王重荣与田令孜的冲突、僖宗出奔及盐池问题见 LT-885-067，核心材料为《旧唐书·僖宗本纪》《旧唐书·王重荣传》与《资治通鉴》光启元年条。冲突责任在传记中高度党派化，盐利的具体收支和地方劳动者的负担并无可连续核算的材料。}

在唐廷尚存的名义之下，这种转变已经在不同区域走出不同道路。王建在891年控制成都，西川由曾为皇帝避难提供条件的军府，变成可独立经营的巴蜀政治中心；杨行密在902年接受吴王封号，则把淮南的实际军政置入唐朝的册命语言之内。前者不等于巴蜀从此与中原隔绝，后者也不等于唐廷仍能调动淮南军队。它们共同表明，节度使名号、王号和诏书仍有价值，却越来越像连接地方军政集团、地方官僚和外部对手的制度接口，而非中央直接统治的证明。%
\footnote{王建控制成都见 LT-891-070；杨行密受封吴王及册封的限度见 LT-902-076。分别可参《旧唐书·王建传》《新五代史·前蜀世家》与《旧唐书·杨行密传》《新五代史·吴世家第一》、相关《资治通鉴》条。两组材料多由后来建立的政权或中原史家保存，成都、淮南各州的服从程度不能由一个称号倒推。}

\section{唐号终结，册命语言没有终结}

晚唐最后的宫廷斗争，把军府化的政治带进了宫门。900年刘季述废昭宗，901年昭宗复位；903年朱温诛宦官并重编神策军，904年又迫昭宗迁洛阳、弑之而立哀帝，次年在白马驿清除反对受禅的朝臣。几次政变的细节在不同史书中并不一致，尤其执行者、诏书和宫禁内部的分工，常由胜利者或清算者叙述。但一条较稳的线索十分清楚：禁军、迁都、朝臣队伍和受禅仪式不再分别属于一个能够自我维持的唐廷；控制洛阳与军队的朱温，能够逐步把这些环节接到自己的宣武军政网络上。%
\footnote{刘季述废立与昭宗复位见 LT-900-074、LT-901-075；朱温诛宦官、迁都弑帝与白马驿清洗见 LT-903-077、LT-904-078、LT-905-079。可互校《旧唐书·昭宗本纪》及刘季述、孙德昭、朱全忠、裴枢诸传与《资治通鉴》唐纪。旧唐书成于后晋，离唐亡较近而保留了较早层次；其仓促编纂和后晋朝廷视角仍要求将宫廷细节同北宋《新唐书》的重整叙事分开阅读。}

907年哀帝禅位，后梁建国，朱温以长期经营的汴州为开封府。受禅的礼文提供了政权更替的形式，却不能消去此前的军事胁迫；开封成为中枢，首先因为那里已有宣武的兵粮、亲军、官僚和汴河交通，而不是因为一个新国号天然拥有全国。由此看五代的“中原王朝”，应理解为能暂时占有开封及其调度网络的政治中心；它所面对的是河东、河北、关中、江淮、巴蜀与岭南等各有兵粮条件的区域，而非一张等待重新上色的空白地图。%
\footnote{唐亡与后梁建国、开封定都见 LT-907-080、LT-907-081，主要依据《旧唐书·哀帝纪》《旧五代史·梁太祖纪一》《新五代史·梁本纪第一》及《资治通鉴》后梁纪。仪式记载能够确定受禅与改制的序列，不能把哀帝意愿或后梁在河南以外的实际控制写成同样确定的事实。}

同一年，各地作出的选择并不相同。王建在成都称帝为前蜀；李茂贞据凤翔而继续使用唐天祐年号；杨渥在淮南拒绝后梁册命；钱镠则接受吴越王号，以杭州和两浙的港口、漕运与地方官僚维持区域政权。它们没有一条共同的“独立时间表”。保留唐年号可以是与后梁相区隔的政治语言，接受后梁王号可以是进行外交和整合地方的资源；称帝也不自动等于对每个州县、山地和港口拥有同样密度的控制。唐亡以后并非只有承认或反对后梁两种选择，而是出现了多种把旧朝礼仪、军府实力和区域资源接合起来的办法。%
\footnote{前蜀建国、岐继续奉唐、吴拒梁命与吴越受封见 LT-907-082、LT-907-083、LT-907-084、LT-907-085。可参《旧五代史·王建传》《旧五代史·李茂贞传》《旧五代史·杨行密传》《旧五代史·钱镠传》，并以《新五代史》前蜀、岐、吴、吴越诸世家和《吴越备史》对读。后出的国别史与北宋史书各有正统立场，王号、年号和都城记录不能替代各政权的赋税账册或基层行政档案。}

福州、广州的情形进一步说明，所谓“十国”不是十个大小相等、制度相同的国家。909年王审知受封闽王，福州军镇借王号联结港口、内陆州县和后梁的名义外交；刘隐死后，刘龑继领岭南军政，917年在番禺称帝，先号大越，旋改称汉。这里可以确认受封、继承与称帝的序列，却不能把后梁的册封读成福建已被直接统治，也不能把南汉的帝号读成岭南全境立即以同一方式纳入广州的财政。港口、山地、州城和军府所能提供的资源并不相同，正是这些差异使区域政权能够存续，也限制了它们的动员深度。%
\footnote{闽王受封见 LT-909-086；刘隐卒、刘龑继领与其称帝改国号见 LT-911-087、LT-917-090。核心材料包括《旧五代史·王审知传》《新五代史·闽世家》《旧五代史·刘隐传》《旧五代史·刘龑传》及《新五代史·南汉世家》。闽、南汉国史多有散佚，现存中原正史和后出国别史保存不均，不能据王号或帝号推断山区、沿海居民的政治归属。}

\section{地方材料所能照亮的，不是一张完整的财政地图}

唐末以前的河西经验，已使这种“名义与实际不同步”的问题变得可见。张议潮在848年驱逐沙州吐蕃守军，唐廷到851年才授其归义军节度使。授节使唐朝得以和河西地方势力建立制度联系，却不表示朝廷已能沿漫长道路恢复常规州县直辖。对于五代十国同样如此：一份册命可以确证朝廷愿意承认谁，一方墓志可以钉住某人的任官、籍贯或死亡日期，一枚钱币或一处城址可以提示某种区域网络；它们都不能单独画出一个政权的固定疆界，更不能把地方精英留下的材料替换成军户、商人、农民和妇女的共同经验。%
\footnote{张议潮起事及唐廷授节见 LT-848-051、LT-851-052，依据《旧唐书·张议潮传》《新唐书·张议潮传》《唐会要》与《资治通鉴》唐纪。张议潮所称“十一州”本身带有奏表和册命的政治性，实际收复范围须逐州检验。}

敦煌文书保存的户籍、计帐、契约和军政文书，能够使研究者追问赋役、寺院经济与地方军府怎样落到纸面；墓志、碑刻和墓券又能为迁徙、任官及政权转换中的个人经历提供独立锚点。可是，敦煌的保存条件特殊，精英墓志也有自我表述与追赠官的规则，二者都不能直接补成开封、成都、杭州、福州或广州的完整财政账。五代十国的国别史又多有散佚：这正是为什么区域材料应当用来打断中原正史过于整齐的叙述，而不能被用来制造另一种同样整齐的地方全史。

《旧唐书》与《新唐书》在这一问题上提供的是两种不同的中原视线。前者成于后晋，距离唐亡较近，保存唐末纪传的层次较早；后者由北宋重修，擅于用方镇表、谱系和制度志把纷乱的关系归入可追溯的格局。两者都不可或缺，也都容易让读者从京师与官僚的成败倒看地方。进入五代后，《旧五代史》承接较近的国史传统，却以中原五朝为骨架，现存版本还有辑佚和复原的层次；《新五代史》叙事有力，十国世家尤便于检索，却带着欧阳修的褒贬与北宋统一后的尺度。因而，后梁、后唐、后晋、后汉、后周的年序不能替代南方诸国的历史，十国的国号也不能被“僭伪”一类后起标签提前裁决。

\section{开封的更替与北方边界的价格}

中原的五朝更迭同样不是一串脱离财政和边防的皇位名单。912年朱友珪弑朱晃自立，显示后梁皇室未能稳住禁军和养子网络；916年阿保机建契丹国并称帝，则使北方出现一个能够以自身制度、部落联盟、汉人官僚和农业区资源进行议价的政治中心。两件事的意义不在于以“内乱”或“外患”两个词盖过具体条件，而在于此后开封的任何继承危机，都必须同时考虑军队对赏赐和指挥的反应，以及北方政权对河东、幽云和交通线的压力。%
\footnote{后梁继承危机见 LT-912-088；阿保机称帝见 LT-916-089。前者主要见《旧五代史·梁太祖纪四》《新五代史·梁本纪第二》，后者可对读《辽史·太祖纪上》《旧五代史·契丹传》与《资治通鉴》后梁纪。朱晃遇弑的宫禁细节和契丹制度的实际运行都不能只按后出中原叙事理解。}

923年李存勖在魏州称帝建后唐，旋即攻入汴州灭后梁；两年后，后唐分道入蜀，前蜀灭亡；但926年庄宗又在兵变中被杀，李嗣源即位。后唐的快速扩张不应被读作重新获得了可永久调度的全国资源。魏博、河东与沙陀军队的集结，汴州官僚财赋的接收，秦凤、峡江方向的进军，都让它能够赢得一时的战役；连续战争后的军赏、军纪和地方接管，却使这套连接很难稳定下来。前蜀覆亡后的巴蜀也没有简单回归中原：孟知祥后来在成都称帝为后蜀，正说明征服军所留下的任官、兵粮和地方军府，可以再次成为区域建国的基础。%
\footnote{后唐灭梁、灭前蜀与926年政变见 LT-923-092、LT-925-094、LT-926-095；后蜀建立见 LT-934-098。主要材料为《旧五代史》唐庄宗、明宗诸纪及孟知祥传，《新五代史》唐本纪、后蜀世家和《资治通鉴》后唐纪。征服者记录擅长安排军功年序，蜀地地方社会的损失、降附与接管仍缺乏可连续量化的证据。}

长江以南的政治并未静止等待中原决定。927年杨溥称帝为吴、马殷受封楚王，分别显示淮南与湖南的军政整合；同样在这一时期，荆南以江陵为中心，借南平王号在后唐、吴、楚之间经营渡口、盐粮与长江交通。937年徐知诰受杨溥禅位，改名李昪，建立南唐；945年南唐灭闽，却没有因此把福建变成一块没有差异的行政腹地。这些事件连在一起，不是为了列出南方政权的增减，而是为了看见江淮、湖湘、荆江、闽岭的水陆通道和军府资源，如何让称臣、受封、称帝、征服与地方自主持久地并存。%
\footnote{吴称帝、楚王受封与荆南受封见 LT-927-096、LT-927-097、LT-924-093；南唐建立及其灭闽见 LT-937-100、LT-945-104。可对读《旧五代史》杨溥、马殷、高季兴、李昪诸传，《新五代史》吴、楚、南平、南唐、闽诸世家及《资治通鉴》。杨氏皇权受徐氏制约、闽地接收不均等限制，正使国号和征服不能被直接换算成统一的行政控制。}

936年石敬瑭在契丹援助下建后晋，承诺割让燕云十六州，把军事援助、政治称臣和边防州郡放到同一场交换中。原始割让文书不存，“十六州”的范围以及后世所谓“儿皇帝”的措辞，都必须警惕后出的中原叙事；但燕云转入契丹控制、后晋以此取得建国援助，是不能略去的结构性变化。942年石重贵拒绝对契丹称臣，礼仪称谓遂与幽云前沿、互市和边境战争纠缠在一起。此处的称谓不是虚浮的面子问题：它把谁拥有指挥权、谁可征发河北河东军民、谁须承担长途战争的费用，都推到更紧张的地步。%
\footnote{后晋建立与燕云问题见 LT-936-099；后晋改变对契丹的外交语调见 LT-941-101、LT-942-102。依据《旧五代史·晋高祖纪一》《旧五代史·晋少帝纪一》《辽史·太宗纪上》及《资治通鉴》后晋纪。燕云割让的文书原件未存，双方使节往返与边境交易亦无完整档案，不能用象征性称谓替代复杂的利益关系。}

\section{从战场接收为军政接收}

944年贝州受围，后晋援军在戚城解围，是一次局部战术转折，未能消除燕云带来的战略劣势。两年后，杜重威等前线将领降契丹，契丹军进入开封，后晋灭亡；辽太宗又在947年北撤，刘知远得以从太原建立后汉、南下进入中原。胜负在这里连续得惊人，却不能被写成某一方天然适合统治中原。契丹军远离草原基地时如何获得粮食、税收和地方合作，后汉如何把河东军、沙陀网络与后晋遗臣接到开封，这些都比“入主”或“收复”更接近政权能否持续的问题。%
\footnote{贝州、戚城战事见 LT-944-103；后晋覆亡、辽军北撤及后汉建立见 LT-946-105、LT-947-106、LT-947-107。可对读《旧五代史》晋少帝纪、汉高祖纪，《辽史·太宗纪下》与《资治通鉴》后晋、后汉纪。围城兵数、辽军北撤的单一原因及开封基层秩序均无足以精确复原的连续材料。}

后汉高祖死后，少年刘承祐即位，枢密、侍卫和三司之间的权力分配日益尖锐；951年郭威在澶州被拥立，入开封建后周，同时刘崇在太原建立北汉并求援于辽。由军队拥立而来的新朝，首先要解决的并非抽象的正统问题，而是禁军、辅政大臣、河东军政和北方援军之间能否形成可用的平衡。954年柴荣在高平击败北汉、辽联军，随后整顿禁军，才使后周能够把北方的防御与此后的财政、宗教和南征政策连接起来。%
\footnote{后汉继承危机见 LT-948-108；后周建立、北汉建立及高平之战见 LT-951-109、LT-951-110、LT-954-111。基本史料为《旧五代史》汉隐帝纪、周太祖纪、周世宗纪，《新五代史》汉本纪、周本纪、东汉世家及《资治通鉴》。太后诏书、拥立程度、临阵退却和辽军投入规模皆有叙事争议，不能以戏剧性细节替代可确认的军事政治后果。}

显德二年整顿寺院，显示所谓财政并不限于田赋或漕粮。后周以废并无敕寺院、收其部分资产的方式，试图取得铜料、劳力、户籍与军费资源；诏令可以确认国家的目标，却不能告诉我们各州执行到何种程度，也不能把寺院废并数和还俗人数直接当作全国社会统计。与其把这一政策解释为单一的宗教敌意，不如把它放在长期战事、铸钱需求与对流动人口进行行政登记的困难之中，同时保留地方寺院、僧尼和施主的反应大多未被保存的事实。%
\footnote{后周整顿寺院见 LT-955-112，依据《旧五代史·周世宗纪二》《五代会要》卷二十六〈寺〉与《资治通鉴》后周纪。诏令和会要可说明规范与中央意图，不能量化地方落实、寺院资产损失或个人信仰。}

956年至958年对南唐淮南的战争，则把“接收”写得更具体。后周沿淮河水陆进军，南唐最终称臣并割让江北诸州；后周随后要安置降军、修复城池、重设州县和漕运据点，才可能把战场变为可供给边防的地区。协议中的州数、边界及移交速度在双方叙事中并不完全一致，所以不能把称臣割地说成一夜之间完成的行政事实；然而，淮河至长江间的兵粮、道路、城池与税收登记从此成为后周军政的一部分，也使长江成为更清楚、但仍不安稳的政权边界。%
\footnote{后周南征、南唐称臣割地与江北接收见 LT-956-113、LT-957-114、LT-958-115。可参《旧五代史·周世宗纪三》《旧五代史·李璟传》《新五代史·南唐世家》《五代会要》卷三十〈州郡〉及《资治通鉴》后周纪。州数、边界和地方赋役成效主要是行政名目，仍须以碑志、考古和后续任官材料逐项检验。}

\section{宋以前，不是秩序以前}

959年柴荣北伐，收复瀛、莫等州后因病班师；幽云核心仍在辽控制之下。柴荣死后，七岁的柴宗训即位，禁军主帅赵匡胤与辅政大臣之间的平衡变得脆弱。960年陈桥兵变后，赵匡胤入开封受禅建宋。宋的建立结束了后梁、后唐、后晋、后汉、后周在开封相继更替的序列，却没有在这一刻抹去辽、北汉、南唐、吴越、后蜀、南汉、楚、荆南等政权和区域网络所形成的世界。将960年当作自然的“重新统一”，会反过来把五代十国写成一段等待终结的空白；更贴切的理解是，新朝继承了开封中枢、禁军与未完成的北方、淮南问题，并将在此后的战争与谈判中处理它们。%
\footnote{后周北伐与柴荣去世见 LT-959-116、LT-959-117；陈桥受禅见 LT-960-118。主要依据《旧五代史·周世宗纪四》《旧五代史·恭帝纪》《辽史·穆宗纪上》《宋史·太祖本纪一》及《资治通鉴》后周纪。北伐的实际收复范围、柴荣托孤与陈桥是否预谋，均不能由宋初正统叙事作确定复原。}

\subsection{交叉阅读窗口九：台城的旧梦与多政权的现实}

五代的开封更替须让《旧五代史》《新五代史》与《资治通鉴》彼此校勘：前者较接近中原国史层次，后者在北宋统一的框架中重整人物与褒贬。对燕云、河东和契丹的部分，还须并读《辽史》、契丹汉文碑刻及高丽方面后出的史料；对河西和南方地方生活，则可使用归义军敦煌文书、墓志、钱币与城址。《五代会要》的州郡、寺院材料能追索中央规范，却不能补成开封、成都、杭州、福州和广州各自的税簿。外部与地方材料不是天然更“真实”，它们同样有保存、译写和政治用途的限度。%
\footnote{五代中原年序可对读《旧五代史》《新五代史》《资治通鉴》后梁至后周纪；契丹问题参《辽史》及相关碑刻，高丽史料为后出编纂，宜与同时性和地域性分别处理。归义军文书、墓志、钱币和城址适于提出区域问题，不能推出诸国的固定疆界或全国赋役。}

韦庄《台城》写“江雨霏霏江草齐，六朝如梦鸟空啼”。此诗为晚唐七言绝句，确切作年难定；诗人借六朝旧都的雨景，将王朝更替压缩成“如梦”的历史感。它的文学力量在于把南方城市的持续景观同政权记忆的断续并置，却绝非吴、吴越、闽、南汉或后蜀财政、贸易与居民态度的调查。把它读入907--960年的多政权世界，应当防止两种捷径：既不把江南写成脱离战争的诗意避难所，也不把“六朝如梦”当作区域秩序必然虚无的证据。钱穆所问的中央与地方如何重接、吕思勉所问的社会经济网络如何延续，都须通过册命、契约、墓志、港口与物质遗存逐项检验；诗只给出一条关于记忆如何塑形的线索。

因此，晚唐至五代十国的历史不能用“中央衰亡、地方割据”八字收束。黄巢战争、盐利冲突、迁都受禅与兵变，使军队同粮道、城市、仓储和地方征发结合得更紧；后梁至后周的朝廷不断试图重新占有这些关系；吴、南唐、吴越、闽、南汉、楚、荆南、前后蜀、岐、北汉以及辽等政权，则以各自不同的年号、都城、册命和资源组合回应这一局面。正史与《通鉴》给我们一条必要的编年脊梁，旧、新唐书和旧、新五代史之间的差异又提醒我们，不要把后来的统一尺度倒灌回当时。地方文书、碑志、钱币和遗址能够校正这条脊梁，却不能填平它周围大量无名的军户、转运者、寺院居民、商人和乡村家庭。多政权并立并非秩序的缺席；它是一种以军事供给、区域财政和有限承认为条件、且始终带着材料缺口的秩序。
